\section{High-Level Language Design}
\label{sec:language}
\subsection{Dialects, Operations, and Attributes}
\label{sec:dialects}
A dialect defines reusable operations, their operand signatures, attributes,
and semantics. Attributes configure operations independently of tensor values:
a transpose, for example, carries \code{permutation: index[]}.
Semantic blocks specify element-type behavior (\code{dtype}), coordinate
bindings, legality, and output order (\code{axes}), and scalar element values
(\code{semantics}). Transpose preserves element type and input values while
changing output coordinate order; \code{yield input} therefore need not mean
an unchanged tensor interface. Appendix~\ref{app:language} gives expanded syntax.

\subsection{Imports and Rewrite Rules}
Rules import dialects with \code{from "tensor.tepl" import TensorLang as t;}
and select operations with \code{use t::{transpose};}.
An LHS pattern precedes its RHS replacement. Nested operations describe
producer--consumer edges, and repeated tensor captures identify the same
matched value. This application graph is distinct from the semantic graph
used to reason about coordinates and values. Our running rule replaces
\code{(transpose (transpose X))} with \code{(transpose X)}.

\subsection{Attribute Binding and Reuse}
\label{sec:binding}
\code{op[@name]} captures an LHS operation's descriptor. Omitted LHS
annotations capture descriptors independently and anonymously; names make
them accessible to conditions, derivations, and compatible RHS operations.
Field access, such as \code{@outer.permutation}, selects one attribute.
\begin{figuretext}{figures/Attribute Binding and Derivation.pdf}{Captured operation descriptors can be read and transformed to configure replacement operations.}{fig:attribute-binding}{An input graph with transpose and dot is replaced by a dot. Colored descriptor
bindings connect the rule to a table with two captured rows and one derived row.}
Reusing \code{@name} on a compatible RHS operation copies its descriptor;
compatibility checks structure, not semantic correctness. A derived descriptor
can configure an operation in the same way.

Figure~\ref{fig:attribute-binding} illustrates a separate transpose-absorption
instance: captured \code{@tr} and \code{@d} feed a user-written
\code{f_transform} that produces \code{@new_d}. Absorbing permutation
\code{[1,0,2]} changes the left contraction position from 1 to 0.
The short fields \code{perm}, \code{l_dim}, and \code{r_dim} are schematic,
rather than the generalized-dot schema used later. The figure shows descriptor
dataflow; arbitrary absorption still requires semantic checking.
\end{figuretext}
\FloatBarrier
\subsection{Writing Attribute Changes}
Metadata expressions read captured attributes and available shape and dtype
bindings without executing tensor computations. Lists, indexing, local
bindings, conditionals, comprehensions, and builtins construct immutable
descriptors. Runtime expressibility is broader than semantic checkability:
an executable expression need not be one the checker can establish correct.

\subsection{Conditions with \texorpdfstring{\code{where}}{where}}
\label{sec:where}
A \code{where} block contains optional \code{let} computations followed by
its applicability condition; only the bindings end with semicolons.
Conditions combine predicates with \code{&&} and \code{||}, and can express
semantic requirements or application policy. TEPL does not infer profitability.
For example, \code{@outer.lhs_batching_dimensions == @inner->B} constrains
axis routing, as explained in Section~\ref{sec:partitioning}.

\subsection{Attribute Construction with \texorpdfstring{\code{derive}}{derive}}
A \code{derive} block computes named RHS descriptors. Composing inner
permutation $p$ with outer permutation $q$ requires $r[i]=p[q[i]]$:
\begin{codeblock}[H]
\begin{lstlisting}[language=TEPL]
rule compose_transposes_explicit {
  (transpose[@outer] (transpose[@inner] X))
  => (transpose[@combined] X)
  derive {
    @combined = {
      permutation: gather(@inner.permutation, @outer.permutation)
    };
  }
}
\end{lstlisting}
\end{codeblock}
\code{where} decides applicability, \code{derive} constructs descriptors, and
the RHS specifies the replacement. Separate local scopes prevent accidental
sharing of temporary bindings; derivations can use captures and earlier
derived descriptors. The explicit constructor is within the supported
transpose-verification fragment; attribute typing alone is insufficient.

\subsection{Automatic Attributes with \texorpdfstring{\code{auto}}{auto}}
\label{sec:auto}
Omitting an RHS descriptor, or writing \code{[auto]}, requests inference:
\begin{codeblock}[H]
\begin{lstlisting}[language=TEPL]
rule compose_transposes_auto {
  (transpose (transpose X)) => (transpose X)
}
\end{lstlisting}
\end{codeblock}
\begin{figuretext}{figures/Introducing auto.pdf}{\code{auto} becomes a reusable attribute function at compilation and is evaluated on matched metadata at runtime.}{fig:introducing-auto}{A vertical flow separates the user rule and operation semantics, compiler
synthesis and verification, and runtime application of the generated function.}
For this rule, compilation derives the same reusable permutation program
$r[i]=p[q[i]]$. Each omitted descriptor is inferred independently;
\code{[@name]} reuses a descriptor, and \code{[auto, field = value]} pins a
field during inference. Users still supply semantics and required conditions.
Unsupported inference produces a diagnostic.

Figure~\ref{fig:introducing-auto} illustrates the compilation/runtime boundary
using matrix transpose absorption. The compiler constructs and checks $f(p,q)$;
each match supplies its transpose permutation and dot descriptor. Absorbing
\code{[1,0]} changes the displayed left contraction position from 0 to 1.
The figure shows only that changing field, not a complete dot definition.
Construction happens once during compilation; runtime evaluates the program
on matched metadata.
\end{figuretext}
\FloatBarrier
\subsection{From \texorpdfstring{\code{auto}}{auto} to Attribute Synthesis}
The compiler must turn known LHS metadata into reusable programs for missing
RHS attributes across ranks and shapes, preserving the computation under the
declared conditions. We now state that requirement precisely.
